\FloatBarrier
\subsection{Error propagation and loss decomposition at the output softmax}
\label{sec:analysis-softmax}

\Cref{sec:analysis-mlp} bounds the perturbation that low-precision arithmetic
injects into one layer.  At the network output, the error travels to the logits
and alters each token's cross-entropy loss.
This subsection analyzes that interface to identify, term by term,
why SR and RN affect the language-model head differently.

\paragraph{Exact decomposition of the cross-entropy loss change.}
Fix one token with reference logits $\lambda\in\R^V$, target class~$y$, and
predictive distribution $p=\softmax(\lambda)$.  Its cross-entropy loss is
\begin{equation}
  \mathcal L(\lambda)=-\lambda_y+\log\sum_{j=1}^Ve^{\lambda_j}.
  \label{eq:token-cross-entropy}
\end{equation}
Let the low-precision run perturb logit $j$ by~$\Delta_j$.  Substituting
$\lambda+\Delta$ into \cref{eq:token-cross-entropy} and subtracting gives
the exact signed loss change
\begin{equation*}
  \delta\mathcal L(\Delta)
  :=\mathcal L(\lambda+\Delta)-\mathcal L(\lambda)
  =-\Delta_y+\log\sum_{j=1}^Vp_je^{\Delta_j}.
\end{equation*}
Let $\widetilde p=\softmax(\lambda+\Delta)$ and let
$D_{\mathrm{KL}}(p\,\|\,\widetilde p)$ denote the Kullback--Leibler (KL)
divergence from $p$ to $\widetilde p$. Direct substitution gives
\begin{equation*}
  D_{\mathrm{KL}}(p\,\|\,\widetilde p)
  =\sum_{j=1}^V p_j\log\frac{p_j}{\widetilde p_j}
  =\log\sum_{j=1}^V p_j e^{\Delta_j}-p^\top\Delta,
\end{equation*}
and therefore, with $g:=p-e_y$, the exact decomposition
\begin{equation}
  \delta\mathcal L(\Delta)
  =g^\top\Delta+D_{\mathrm{KL}}(p\,\|\,\widetilde p).
  \label{eq:one-token-kl-decomposition}
\end{equation}
The linear term is signed and records the target-specific shift, whereas the
KL term is non-negative and records the change in the full predictive
distribution.

\paragraph{Local quadratic form of the KL term.}
The KL divergence satisfies $D_{\mathrm{KL}}(p\,\|\,p)=0$ and has vanishing gradient
at $\Delta=0$. Its local quadratic approximation is governed by the Fisher
information matrix of the categorical distribution parameterized by logits,
\begin{equation*}
  H = \diag(p) - pp^\top,
\end{equation*}
which coincides with the Hessian of the cross-entropy loss $\nabla^2_\lambda \mathcal L(\lambda)$
\citep[Section~6]{martens2020new}. Although the KL divergence is asymmetric, both
orientations share the same second-order expansion at $\Delta=0$:
\begin{equation}
  D_{\mathrm{KL}}(p\,\|\,\widetilde p)
  =\tfrac12\,\Delta^\top H\Delta+r,
  \qquad r=\mathcal O(\|\Delta\|^3).
  \label{eq:logit-kl-second-order}
\end{equation}
Because $H$ is the covariance matrix of a single categorical draw with probabilities $p$,
for any perturbation $\Delta \in \R^V$ the quadratic form equals the variance across classes:
\begin{equation}
  \Delta^\top H \Delta = \sum_{j=1}^V p_j \Delta_j^2 - (p^\top\Delta)^2 = \Var_p(\Delta) \ge 0,
  \label{eq:fisher-variance-identity}
\end{equation}
confirming that $H \succeq 0$ and that the local KL penalty charges variations in
$\Delta$ across classes considered plausible by $p$.
Fisher information similarly guides parameter sensitivity in post-training
quantization (e.g., SqueezeLLM, \citealp{kim2024squeezellm}); in logit space,
$H$ arises analytically to penalize non-uniform operational perturbations directly.

An exact consequence of the logit formulation is \emph{shift invariance}
\citep{blanchard2021accurately}: any uniform perturbation $\Delta = c\bfone$ leaves
the predictive distribution identical ($\widetilde p = \softmax(\lambda + c\bfone) = p$),
so that $D_{\mathrm{KL}}(p\,\|\,\widetilde p) = 0$.
Likewise, $g^\top (c\bfone) = c(p-e_y)^\top\bfone = 0$, ensuring that $\delta\mathcal L(c\bfone) = 0$
exactly. The one-dimensional subspace $\operatorname{span}(\bfone)$ constitutes the null
space of the Fisher metric ($H\bfone = \mathbf{0}$).

\paragraph{The two probability spaces.}
Conditional on a fixed token~$w$, two distinct probability spaces arise for SR
and RN.
\begin{enumerate}[leftmargin=1.7em,label=(\roman*)]
  \item Under SR, every rounded elementary operation along the trajectory draws
        fresh pseudorandom bits. This operational randomness makes the
        logit perturbation $\Delta_{\SR}^{(w)}$ a random vector, driven by
        elementary rounding errors that satisfy the mean-independence of \cref{sec:analysis-mlp}.

  \item Under RN, every operation returns one deterministic value: there is no
        operational randomness. Conditional on a fixed token,
        \begin{equation}
          \E_{\RN}[\Delta_{\RN}^{(w)}\mid w]=\Delta_{\RN}^{(w)},
          \qquad
          \Sigma_{\RN}^{(w)}:=\operatorname{Cov}_{\RN}(\Delta_{\RN}^{(w)}\mid w)=0.
          \label{eq:rn-degenerate}
        \end{equation}
\end{enumerate}

\paragraph{Drift--variance decomposition.}
For either mode $m\in\{\SR,\RN\}$, decompose the logit perturbation as
\begin{equation}
  \Delta_m=b_m+\xi_m,
  \qquad
  \E_m[\xi_m]=0,
  \label{eq:local-bias-cov}
\end{equation}
where $b_m:=\E_m[\Delta_m]$ is the drift and $\xi_m$ is the centered
fluctuation, with covariance $\Sigma_m:=\operatorname{Cov}_m(\xi_m)$.
By \cref{eq:rn-degenerate}, on a fixed token RN carries
$b_{\RN}=\Delta_{\RN}$ and $\Sigma_{\RN}=0$: its entire perturbation is
drift.

Substituting $\Delta_m=b_m+\xi_m$ into the local KL expansion
\cref{eq:logit-kl-second-order} and taking expectations $\E_m[\cdot]$ over
operational randomness (over rounding draws for SR, and degenerate for
deterministic RN), the linear term in \cref{eq:one-token-kl-decomposition} yields
$\E_m[g^\top\Delta_m]=g^\top b_m$.
The quadratic term expands as
\begin{equation*}
  \Delta_m^\top H\Delta_m
  =b_m^\top Hb_m + 2b_m^\top H\xi_m + \xi_m^\top H\xi_m.
\end{equation*}
Taking expectations, the cross term vanishes since $\E_m[\xi_m]=0$, and the
quadratic fluctuation term evaluates to $\E_m[\xi_m^\top H\xi_m] = \operatorname{tr}(H\Sigma_m)$.
Thus the KL contribution itself separates as
\begin{equation*}
  \E_m[D_{\mathrm{KL}}(p\,\|\,\widetilde p_m)]
  =\tfrac12\,b_m^\top Hb_m +\tfrac12\operatorname{tr}(H\Sigma_m) +R_m,
\end{equation*}
where $R_m:=\E_m[r]$.  Combining this KL decomposition with the linear
target-shift term yields the expected excess loss:
\begin{equation}
  \E_m[\delta\mathcal L]
  =\underbrace{g^\top b_m}_{D_m:\text{ drift}}
  +\underbrace{\tfrac12\,b_m^\top Hb_m}_{Q_m:\text{ curvature of drift}}
  +\underbrace{\tfrac12\operatorname{tr}(H\Sigma_m)}_{V_m:\text{ variance penalty}}
  +\underbrace{R_m}_{\text{remainder}},
  \label{eq:head-full-comparison}
\end{equation}
where $D_m$ is signed; $Q_m$ and $V_m$ are non-negative by $H\succeq0$.
Removing drift is not free: an unbiased mode avoids the drift terms $D_m+Q_m$,
but pays the variance penalty $V_m$.

\subsection{Empirical validation of the drift--variance trade-off}
\label{sec:analysis-empirical-drift}

We evaluate this trade-off on DistilGPT-2 using $N=1020$ scored WikiText-2
positions, with reference loss $\mathcal L_0=4.1354$ nats.
The decomposition uses $32$ SR seeds at $t=4,\ldots,10$, with one
deterministic RN run per precision. We use bias-corrected quadratic
estimators.  The logit measurement protocol and the estimators are detailed in \cref{app:drift-measurements}.

\paragraph{SR avoids systematic drift but incurs local fluctuation.}
The decomposition $\Delta_m = b_m + \xi_m$ (\cref{eq:local-bias-cov})
splits each logit perturbation into drift $b_m = \E_m[\Delta_m]$ and centered
fluctuation $\xi_m$. For RN the split is trivial: $b_{\RN} = \Delta_{\RN}$
and $\xi_{\RN} = 0$, so the perturbation at each token--coordinate pair is
entirely deterministic drift. For SR, the perturbation is stochastic,
so its expected magnitude is measured by its operational root mean square.
All plotted quantities represent the operational statistics of a
single forward pass; multiple seeds ($S=32$) are used exclusively to construct
empirical estimators for these quantities.

\Cref{fig:drift-distribution} visualizes the three quantities that enter the
loss decomposition (\cref{eq:head-full-comparison}). Row~(i) shows the
operational perturbation magnitude per run across all token--coordinate pairs.
At the MLP down-projection, RN accumulates deterministic drift, whereas SR's
per-run perturbation is smaller. At the head, however, SR's per-run perturbation
is larger than RN's. This does not contradict the worst-case bounds of
\cref{thm:vector-bc}. At the head, inputs are centered by LayerNorm;
empirically, we observe that this centering induces cancellation in RN rounding
errors, preventing the accumulation of systematic drift. Because the projection
is linear, SR guarantees unbiasedness, but incurs higher local variance to
do so (local errors up to $1.0u$, compared to RN's strict $0.5u$ bound). In the
absence of systematic drift, the accumulated perturbation is governed by this
local variance, which is consistent with SR perturbations being larger than RN
at the head.

Nevertheless, not all of this perturbation magnitude translates to excess loss.
By shift invariance (\cref{eq:fisher-variance-identity}), a per-token uniform
shift leaves softmax probabilities unchanged, so only the non-uniform
component feeds into $Q_m$. Row~(ii) isolates this loss-relevant component
by centering the drift (subtracting each token's vocabulary mean from $b_j$).
The reduction is pronounced at the MLP: RN's spread shrinks by an order
of magnitude, confirming that most of its accumulation was uniform.
SR's centered drift is near-zero everywhere.

Row~(iii) removes the uniform component of the fluctuation (to which the
softmax is insensitive by shift invariance). Let $\tilde\xi_j = \xi_j - p^\top\xi$
denote the $p$-centered operational fluctuation; the row reports its standard
deviation $\mathrm{std}(\tilde\xi_j)$. RN is deterministic
($\xi_{\RN}=0$). The centered fluctuation is smaller at the MLP than at
the head.
MLP noise is predominantly a uniform shift of the logits, whereas
head noise changes the differences between logits and therefore the softmax
output.
\begin{figure}[!htb]
  \centering
  \safeincludegraphics[width=\linewidth]{figures/drift_boxplot.pdf}
  \vspace{-1em}
  \caption{Operational decomposition of logit perturbations across all $51{,}262{,}140$ token--coordinate pairs (for the MLP, rounding errors are injected at the down-projection and propagated to the final logits).
  \textbf{(i)} Per-run RMS perturbation for pair $j$: $|\Delta_j|$ for deterministic RN ($\xi_j=0$), and $\sqrt{\mathbb{E}[\Delta_j^2]} = \sqrt{b_j^2 + \operatorname{std}(\xi_j)^2}$ for SR.
  \textbf{(ii)} Centered drift $\tilde b_j$ (per-token vocabulary mean removed): non-uniform component charged via $Q_m$.
  \textbf{(iii)} Operational standard deviation of the $p$-centered fluctuation $\mathrm{std}(\tilde\xi_j)$, where $\tilde\xi_j = \xi_j - p^\top\xi$; RN: $\xi_{\RN}=0$.
  Boxes IQR, median line, 1st--99th-percentile whiskers; vertical scales differ.}
  \label{fig:drift-distribution}
\end{figure}

\paragraph{The role of input distributions.}
The accumulation of deterministic rounding errors in inner products depends heavily on the distribution of the operands \citep{elarar2022error}. When inputs are predominantly positive (such as the outputs of the GELU activation at the MLP down-projection), RN rounding errors become biased and, depending on the weight distribution, systematically accumulate into a large drift. Conversely, when inputs are zero-centered (such as the outputs of LayerNorm at the language-model head), positive and negative RN errors tend to cancel, preventing the accumulation of systematic drift.

\paragraph{Loss decomposition at the language-model head.}
\label{par:head-cost-unbiasedness}
Although SR produces smaller logit drift at the head, it incurs higher excess
loss than RN throughout $t=4,\ldots,10$
(\cref{tab:decomposition,fig:head-decomposition-terms}). The decomposition in
\cref{eq:head-full-comparison} explains the reversal. RN is deterministic
($V_{\RN}=0$), so its perturbation contributes only through $D_{\RN}$,
$Q_{\RN}$, and $R_{\RN}$. SR's operational randomness introduces non-uniform
fluctuations---the large $p$-centered fluctuation visible in
\cref{fig:drift-distribution}~(iii)---that incur the quadratic cost
$V_{\SR}=\tfrac12\operatorname{tr}(H\Sigma_{\SR})$. At the head the noise is
genuinely non-uniform across the vocabulary, so the softmax cannot ignore it.

Over $t=6,\ldots,9$, the variance penalty $V_{\SR}$ accounts for $100$--$105\%$ of SR's
excess loss, while $D_{\SR}$ and $Q_{\SR}$ lie within $1.6$ standard errors of
zero. For RN, the perturbation carries no fluctuation ($V_{\RN}=0$);
its non-uniform drift produces $Q_{\RN}=1.58$ and $R_{\RN}=0.50$
nats at $t=6$, for a total excess of $2.12$ nats---less than half of SR's
$4.40$ nats.

\paragraph{Loss decomposition at the MLP down-projection.}
At the MLP down-projection, the loss ordering reverses:
SR has an excess loss of $0.05$ nats against RN's $0.24$ at $t=6$
(\cref{fig:mlp-decomposition-terms,tab:decomposition}).
The 3072-length inner product allows deterministic rounding errors under RN
to align across many terms, producing a wide, non-uniform drift dispersion
that incurs a curvature penalty of $Q_{\RN}=0.29$ nats.

SR prevents this systematic alignment, keeping $Q_{\SR}$ at $0.017$ nats.
Although SR's raw fluctuation is comparable to or larger than at the head
($6.23$ versus $4.27$ logits RMS at $t=6$), the noise has far fewer degrees of
freedom. At the head, SR injects rounding errors directly into all $50{,}257$
logits, scrambling the differences between them. At the MLP down-projection, SR
injects noise into the $768$-dimensional hidden state. Restricted to these $768$
degrees of freedom, the $50{,}257$ logits are highly coupled.  This
dimensionality restriction likely explains why the MLP fluctuation manifests
predominantly as a uniform shift across the vocabulary, which the softmax simply
ignores. The near-uniform shift also depends on the trained output weights and
the directions of the propagated error. As a result, the $p$-centered
fluctuation in \cref{fig:drift-distribution}(iii) collapses at the MLP relative
to the head, yielding $V_{\SR}=0.083$ nats, fifty times smaller than at the
head. The reduction in drift curvature therefore more than compensates for SR's
variance penalty.

\begin{figure}[!htb]
  \centering
  \safeincludegraphics[width=\linewidth]{figures/site_decomposition_terms.pdf}
  \caption{Stacked decomposition of expected excess loss
    $\E_m[\delta\mathcal L] = D_m + Q_m + V_m + R_m$ at the head and MLP
    down-projection at $t=6,7,8$. Positive components stack upward from zero
    and negative components downward; solid horizontal marks indicate the
    measured mean excess loss (with one standard error for SR; RN is
    deterministic). At the head, the variance penalty $V$
    dominates SR excess loss; at the MLP down-projection, RN is dominated by drift
    curvature $Q$. Full values are in \cref{tab:decomposition}.}
  \label{fig:head-decomposition-terms}
  \label{fig:mlp-decomposition-terms}
\end{figure}

